\exercisesection{Smooth algebras}

\begin{enumerate}
\def\labelenumi{\arabic{enumi})}
\tightlist
\item
  Let \(k\) a ring, A a linearly topologized formally smooth \(k\) -linearly topologized algebra. One says that A is a \(k\) -algebra formally étale (resp. formally net) if it satisfies the following condition: whatever the \(k\) -algebra C, endowed with the discrete topology, and the ideal N of square zero of C, every continuous homomorphism of \(k\) -algebras A → C/N admits a unique (resp. at most one) lifting A → C. For A to be formally étale over \(k\) it is necessary and sufficient that it be formally smooth and formally unramified. Every quotient of \(k\) is formally unramified.
\end{enumerate}

\begin{enumerate}
\def\labelenumi{\alph{enumi})}
\item
  Prove the analogues of the statements of Props. 3 and 4 for formally smooth or formally étale algebras.
\item
  Let J be an ideal of \(\Lambda\) such that the topology of A is the J-adic topology; denote by \(\widehat{\Omega}_{k}(\Lambda)\) the separated completion of the A-module \(\Omega_{k}(A)\) (for the J-adic topology). For the \(k\) -topological algebra A to be formally smooth, it is necessary and sufficient that \(\widehat{\Omega}_{k}(A)\) be zero (observe that each of these properties is equivalent to the fact that every \(k\) -derivation of A into an A-module annihilated by a power of J is zero).
\end{enumerate}

\begin{enumerate}
\def\labelenumi{\arabic{enumi})}
\setcounter{enumi}{1}
\item
  Let A be a Noetherian ring, J an ideal of A, and \(\widehat{\mathbf{A}}\) the separated completion of A for the J-adic topology. If the \(\Lambda\) -algebra \(\widehat{\mathbf{A}}\) is formally smooth (for the discrete topology), it is formally étale (observe that one has \(\Omega_{\mathbf{A}}(\widehat{\mathbf{A}}) = J^{n}\Omega_{\mathbf{A}}(\widehat{\mathbf{A}})\) for every integer \(n\) ).
\item
  Let K be a finitely generated extension of a field \(k\) ; prove that the following conditions are equivalent:
\end{enumerate}

\begin{enumerate}
\def\labelenumi{(\roman{enumi})}
\item
  the k-algebra K is formally smooth;
\item
  the \(k\) -algebra K is formally étale;
\item
  K is a finite separable extension of k.
\end{enumerate}

\begin{enumerate}
\def\labelenumi{\arabic{enumi})}
\setcounter{enumi}{3}
\tightlist
\item
  Let \(k\) be a Noetherian ring, A a finitely generated \(k\) -algebra. One says that \(\Lambda\) is an \(k\) -étale algebra (resp. smooth) if it is formally étale (resp. formally smooth) when equipped with the discrete topology.
\end{enumerate}

\begin{enumerate}
\def\labelenumi{\alph{enumi})}
\tightlist
\item
  Prove that the following conditions are equivalent:
\end{enumerate}

\begin{enumerate}
\def\labelenumi{(\roman{enumi})}
\item
  the k-algebra A is smooth;
\item
  \(\Omega_k(A) = 0\)
\item
  the \(\mathrm{A}\otimes_{k}\mathrm{A}\) -module A is projective.
\end{enumerate}

(Let I be the kernel of the canonical homomorphism \(\mathrm{A} \otimes_k \mathrm{A} \to \mathrm{A}\) ; show that (ii) and (iii) are both equivalent to

\begin{enumerate}
\def\labelenumi{(\roman{enumi})}
\setcounter{enumi}{3}
\tightlist
\item
  \(I_{q} = 0\) for every prime ideal q of A \(\otimes_{k}\) A containing I.)
\end{enumerate}

\begin{enumerate}
\def\labelenumi{\alph{enumi})}
\setcounter{enumi}{1}
\item
  Let \(f \in k[T]\) be a polynomial such that f and \(f'\) generate the unit ideal of k{[}T{]}. Prove that the k-algebra \(k[T]/(f)\) is étale (cf.~no. 3, Example 3).
\item
  If k is a field, the smooth k-algebras are the étale k-algebras, and this notion coincides with that introduced in A, V, p.~28 (the case of algebras of finite degree over k follows from loc. cit., p.~32, th. 3; if A is a smooth k-algebra, consider a maximal ideal m of A and apply the preceding case to A/m \(^{n}\) ). The smooth k-algebras are therefore the finite products of finite separable extensions of k (A, V, p.~34, th. 4).
\item
  For the k-algebra A to be smooth, it is necessary and sufficient that for every prime ideal p of k the \(\kappa(\mathfrak{p})\) -algebra \(\kappa(\mathfrak{p})\otimes_{k}A\) be étale.
\end{enumerate}

\begin{enumerate}
\def\labelenumi{\arabic{enumi})}
\setcounter{enumi}{4}
\tightlist
\item
  Let A be a complete Noetherian local ring and B an essentially finitely generated A-algebra.
\end{enumerate}

\begin{enumerate}
\def\labelenumi{\alph{enumi})}
\item
  Show that the set of prime ideals p of A such that the ring \(\Lambda_{\mathfrak{p}}\) is regular is open in \(\operatorname{Spec}(\mathbf{A})\) (one will use Exercise 16 of VIII, § 5 to reduce to the case where A is integral, then IX, § 2, no. 5, th. 3 and Exercise 17 of VIII, § 5).
\item
  Suppose that A contains a field \(^{1}\) . Show that the set of prime ideals q of B such that the ring \(B_{q}\) is regular is open in Spec(B) (one will use § 7, no. 9, cor. 2 of th. 3 and IX, § 3, no. 3, th. 2).
\end{enumerate}

\begin{enumerate}
\def\labelenumi{\arabic{enumi})}
\setcounter{enumi}{5}
\item
  Let \(k\) a ring, A a linearly topologized formally smooth \(k\) -algebra. Prove that for every open ideal J of A, the A/J-module \((\mathrm{A} / \mathrm{J}) \otimes_{\mathrm{A}} \Omega_k(\mathrm{A})\) is projective (prove, using the example of no. 1, that for every surjective homomorphism of A/J-modules \(\mathrm{M} \to \mathrm{M}''\) , every \(k\) -derivation of A into \(\mathrm{M}''\) lifts to M).
\item
  Let \(k\) a ring, \(\rho : A \to B\) a homomorphism of \(k\) -algebras, \(J\) an ideal of \(B\) . One says that \(B\) , equipped with the \(J\) -adic topology, is formally smooth over \(A\) relative to \(k\) if for every \(A\) -algebra \(C\) (which we endow with the discrete topology) and every ideal of square zero \(N\) of \(C\) , every continuous homomorphism of \(A\) -algebras \(B \to C/N\) that lifts to a continuous homomorphism of \(k\) -algebras from \(B\) in \(C\) also lifts to a continuous homomorphism of \(A\) -algebras from \(B\) in \(C\) .
\end{enumerate}

\begin{enumerate}
\def\labelenumi{\alph{enumi})}
\tightlist
\item
  Show that the following conditions are equivalent:
\end{enumerate}

\begin{enumerate}
\def\labelenumi{(\roman{enumi})}
\item
  B is formally smooth over A relative to k for the J-adic topology.
\item
  for all \(k\) -derivation D of \(\Lambda\) In a B-module E annihilated by a power of J, there exists a \(k\) -derivation \(\tilde{D}: B \to E\) such that \(\tilde{D} \circ \rho = D\) ;
\item
  for every integer \(n \geqslant 0\) , the canonical map \(u_n: \Omega_k(A) \otimes_A (B/J^n) \longrightarrow \Omega_k(B) \otimes_B (B/J^n)\) is injective and its image is a direct summand.
\end{enumerate}

\begin{enumerate}
\def\labelenumi{\alph{enumi})}
\setcounter{enumi}{1}
\tightlist
\item
  We assume B is formally smooth over k for the J-adic topology. Prove that for B to be formally smooth over A, it is necessary and sufficient that the canonical homomorphism \(u_{1}:\Omega_{k}(A)\otimes_{A}(B/J)\to\Omega_{k}(B)\otimes_{B}(B/J)\) be injective and that its image be a direct summand (use a) and exerc. 6 to prove that a retraction of \(u_{1}\) lifts to a retraction of \(u_{n}\) for every n).
\end{enumerate}

¶ 8) Let k be a field and T a finite family of indeterminates. One equips the k-algebra \(k[[T]]\) with the discrete topology.

\begin{enumerate}
\def\labelenumi{\alph{enumi})}
\item
  Prove that \(k[[T]]\) is formally smooth over \(k[T]\) with respect to k (exerc. 7).
\item
  If \(k[[\mathbf{T}]]\) is formally smooth over \(k\) , prove that the characteristic \(p\) of \(k\) is strictly positive and that \(k\) is a finite extension of \(k^p\) (deduce from exerc. 2 that \(\Omega_{k|\mathbf{T}|}(k[[\mathbf{T}]])\) is zero, and consequently that \(\Omega_{k(\mathbf{T})}(k((\mathbf{T})))\) is zero; then prove, using prop. 6 of A, V, p.~99, that the subextension of \(k((\mathbf{T}))\) consisting of series whose coefficients generate a finitely generated extension of \(k^p\) is equal to \(k((\mathbf{T})))\) .
\end{enumerate}

\begin{enumerate}
\def\labelenumi{\arabic{enumi})}
\setcounter{enumi}{8}
\tightlist
\item
  Let \(k\) a field of characteristic \(p > 0\) , containing a decreasing filtered family \((k_{\lambda})_{\lambda \in \Lambda}\) of subfields satisfying \(\bigcap_{\lambda \in \Lambda} k_{\lambda} = k^{p}\) .
\end{enumerate}

\begin{enumerate}
\def\labelenumi{\alph{enumi})}
\item
  Prove that the intersection of the kernels of the canonical homomorphisms from \(\Omega(k)\) in \(\Omega_{k_{\lambda}}(k)\) , for \(\lambda\) traversing \(\Lambda\) is reduced to (0) (consider a \(p\) -basis of \(k\) ).
\item
  Let A be a regular local k-algebra, endowed with the topology \(m_{A}\) -adic. Prove that for A to be formally smooth over k, it is necessary and sufficient that \(\Lambda\) let be formally smooth over k relative to \(k_{\lambda}\) , for all \(\lambda \in \Lambda\) (using a), Prop. 8, and Exerc. 7, a)).
\end{enumerate}

¶ 10) Let A be a complete Noetherian local ring, K its field of fractions. Let p be a prime ideal of A; denote by B the local ring \(A_{p}\) and \(\widehat{B}\) its completion.

\begin{enumerate}
\def\labelenumi{\alph{enumi})}
\item
  Assume A is regular and K has characteristic zero. Show that the K-algebra \(K \otimes_{B} \widehat{B}\) is absolutely regular.
\item
  Assume A is regular and K has characteristic p. \textgreater{} 0, so that A is identified with a ring of formal power series. \(k[[T_{1}, \ldots, T_{n}]]\) (IX, § 3, no. 3, th. 2). We denote by \((k_{\lambda})_{\lambda \subset \Lambda}\) the family of subfields \(k_{\lambda}\) of k containing \(k^{p}\) and such that k is a finite-degree extension of \(k_{\lambda}\) . Let \(\lambda \in \Lambda\) ; we denote \(A_{\lambda}\) the ring \(k_{\lambda}[[T_{1}^{p}, \ldots, T_{n}^{p}]]\) , \(B_{\lambda}\) the local ring of \(A_{\lambda}\) into the prime ideal \(p \cap A_{\lambda}\) , and by \(K_{\lambda}\) the field of fractions of \(\Lambda_{\lambda}\) and \(B_{\lambda}\) Prove that the A-algebra B (endowed with the discrete topology) is formally smooth relative to \(B_{\lambda}\) (one may first show that \(\widehat{B}\) is isomorphic to \(B \otimes_{B_{\lambda}} \widehat{B_{\lambda}}\) ).
\item
  Under the hypotheses of b), prove that every local ring of \(K \otimes_{B} \widehat{B}\) is of the form \((\widehat{B})_{\mathbf{r}}\) , with \(\mathfrak{r} \in \operatorname{Spec}(\widehat{B})\) and \(\mathfrak{r} \cap B = \{0\}\) , and that it is formally smooth over K relative to \(K_{\lambda}\) Deduce from this that the K-algebra \(K \otimes_{B} \widehat{B}\) is absolutely regular (prove that the family \(\{K_{\lambda}\}_{\lambda \in \Lambda}\) of subfields of K is downward filtering and that \(\bigcap_{\lambda \in \Lambda} K_{\lambda} = k^{p}((T_{1}, \ldots, T_{n}))\) ; then apply exercise 9).
\item
  We return to the general case. Prove that A is a Grothendieck ring (§ 6, exerc. 5); if \(\mathfrak{q} \in \operatorname{Spec}(B)\) , we need to show that the \(\kappa(\mathfrak{q})\) -algebra \(\kappa(\mathfrak{q}) \otimes_{\mathrm{B}} \widehat{\mathrm{B}}\) is absolutely regular; replacing A by \(\mathrm{A}/(\mathrm{A} \cap \mathfrak{q})\) one may assume \(\Lambda\) integral domain and \(\mathfrak{q} = \{0\}\) ; we will use IX, § 2, n° 3, th. 3 and § 3, n° 3, th. 2 to reduce to the case where A is a complete regular local ring.
\end{enumerate}

\begin{enumerate}
\def\labelenumi{\arabic{enumi})}
\setcounter{enumi}{10}
\item
  Let A be a Noetherian ring; suppose that for every maximal ideal m of A, the completion of the local ring \(A_{m}\) is an \(A_{m}\) absolutely regular algebra. Show that A is a Grothendieck ring (§ 6, exerc. 5). (Use loc. cit. and the preceding exercise to show that for every maximal ideal m of A, the ring \(A_{m}\) is a Grothendieck ring.)
\item
  Let A be a Noetherian local ring. For A to be a Grothendieck ring, it is necessary and sufficient that, for every finite integral A-algebra B, every maximal ideal m of B, and every prime ideal q of the completed ring \(\widehat{B_{m}}\) satisfying \(B_{m} \cap q = (0)\) , the local ring \((\widehat{B_{m}})_{q}\) let be regular (to show that this condition is necessary, it suffices, by the previous exercise, to show that, for every prime ideal p of A, and every finite extension K of \(\kappa(\mathfrak{p})\) , the ring \(K \otimes_{A} \widehat{A}\) is regular; one may introduce a finite integral A-algebra B with fraction field \(K\) to note \(\mathfrak{m}_{1},\ldots,\mathfrak{m}_{r}\) its maximal ideals, and interpret the local rings of \(K\otimes_{A}\widehat{A}\) as local rings of one of the \(\widehat{B_{\mathfrak{m}_{i}}}\) .
\item
  Let \(k\) a field and A a \(k\) -algebra essentially of finite type. Show that A is a Grothendieck ring (reduce to the case where A is a local ring of a polynomial ring over \(k\) in a finite number of indeterminates; using the preceding exercise, reduce next to showing that, for every prime ideal \(\mathfrak{r}\) of \(\mathrm{A}[\mathrm{T}_1,\ldots ,\mathrm{T}_n]\) , every local ring C of \(\mathrm{A}[\mathrm{T}_1,\ldots ,\mathrm{T}_n]\) at a maximal ideal containing \(\mathfrak{r}\) , and every ideal \(\mathfrak{q}\) of the completed ring \(\widehat{\mathrm{C}}\) such that \(\mathrm{C}\cap \mathfrak{q} = \mathfrak{r}\mathrm{C}\) , the ring \((\widehat{\mathrm{C}})_q / \mathfrak{r}(\widehat{\mathrm{C}})_q\) is regular. For this, one may use th. 3 of no. 9).
\end{enumerate}

¶ 14) Let A be a local Grothendieck ring. Show that the set of prime ideals p of A such that the ring \(A_{p}\) is regular is open in \(\operatorname{Spec}(A)\) (let \(\widehat{A}\) the completion of A; we will use exercise 5, exercise 16 of VIII, § 5, and the fact that the canonical mapping \(\operatorname{Spec}(\widehat{A}) \to \operatorname{Spec}(A)\) is open).

¶ 15) Let \(k\) a perfect field of characteristic \(p > 0\) , \((\mathrm{T}_{n})_{n \in \mathbf{N}}\) a family of indeterminates, and K the field \(k((\mathrm{T}_{n})_{n \in \mathbf{N}})\) . Let \(\Lambda\) the subring of the ring of formal power series \(\mathrm{K}[[\mathrm{X},\mathrm{Y}]]\) generated by \(\mathrm{K}^{p}[[\mathrm{X},\mathrm{Y}]]\) and K. The ring A is a regular local Noetherian ring of dimension 2, and its completion is \(\mathrm{K}[[\mathrm{X},\mathrm{Y}]]\) For every integer \(n\) , we denote \(p_{n}\) the element \(\mathrm{X} - \mathrm{T}_{n}\mathrm{Y}\) of A, \(q_{n}\) the product \(p_{0} \ldots p_{n}\) , and we set \(c = \sum_{n \in \mathbf{N}} q_{n}\mathrm{T}_{n}\) We denote by B the sub-A-algebra dc \(\mathrm{K}[[\mathrm{X},\mathrm{Y}]]\) generated by \(c\) .

\begin{enumerate}
\def\labelenumi{\alph{enumi})}
\item
  Let \(n\) an integer \(> 0\) and \(\mathfrak{p}\) a prime ideal of \(B\) of height 1 that contains \(p_n\) Show that the ring \(B_{\mathfrak{p}}\) is not normal (for every integer \(m\) , we will set \(c_m = (c - \sum_{i=0}^{m-1} q_i T_i)/q_m\) and we will show that \(B_{\mathfrak{p}} = A_{A p_n}[c_{n-1}]\) , so that \(c_n\) is not in \(B_{\mathfrak{p}}\) , whereas \(c_n^p\) belongs to \(A\) ).
\item
  Deduce that the set of prime ideals p of B such that the ring \(B_{p}\) is regular contains no nonempty open subset of Spec(B) (note that B is a faithfully flat A-module).
\end{enumerate}

